%%% mode: latex
%%% TeX-master: t
%%% End:

\chapter{总结与展望}
\label{cha:conclusion}

\section{本文主要内容及结论}
\label{sec:summary}

医疗大语言模型作为人工智能在医疗领域的重要应用，其安全性直接关系到临床决策的正确性和患者的生命安全。然而，随着多智能体协作系统在医疗场景的广泛应用，智能体之间的信息传递可能放大错误并产生级联效应，对患者安全构成潜在威胁。当前研究主要聚焦于单体模型的安全性评估，缺乏针对多智能体协作场景的系统性安全评估框架。本文针对上述问题提出解决方法，并通过系统性的实验验证所提出方法的有效性。总体而言，本文主要完成的工作内容如下：

（1）本文设计了一个面向医疗多智能体系统的自动化安全评估框架MLSE（Medical LLM Safety Evaluation）。该框架基于AutoGen和Ollama技术栈，采用模块化设计思想，包含场景配置、攻击执行、安全评估和报告生成四个核心模块。框架创新性地引入了Moderator Agent、Target Cascade Agents、Safety Agent、Evaluator Agent、Prompt Mutation Agent和Performance Agent等六类专业化智能体角色，实现了从攻击模拟、安全检测、性能评估到结果聚合的全流程自动化。通过llama.cpp模型量化和Ollama统一接口，框架实现了对BioGPT、AlpaCare、Asclepius、BioMistral、BioLlama3等多个医疗大模型的快速适配和高效部署。

（2）本文系统评估了单体医疗大语言模型的对抗鲁棒性。针对数据投毒攻击，实验结果表明，随着投毒比例从1\%增加到10\%，模型的攻击成功率呈现明显上升趋势。针对提示词操纵攻击，医疗大模型甚至比通用大语言模型Llama2表现出更高的攻击成功率。值得注意的是，传统的文本质量指标如困惑度和连贯性在攻击前后并未发生显著变化，表明攻击具有极强的隐蔽性，常规的质量检测方法难以有效识别。

（3）本文构建了面向临床诊疗链的多智能体级联安全评估子系统，并首次提出了错误放大因子$\alpha$的数学定义。该指标用于量化多智能体协作系统中错误传播的放大效应，为理解级联错误的传播机制提供了定量分析工具。实验设计了16种不同的智能体配置组合，结果揭示了多智能体协作系统中的级联错误传播规律：首个智能体的状态对整个链式协作系统的安全性具有决定性影响，一旦入口智能体被成功攻击，后续智能体即使本身未受攻击也会输出错误结果。实验中计算得到的错误放大因子$\alpha$显著大于1，验证了理论模型中提出的放大假设。

（4）本文构建了基于真实世界类风湿关节炎临床数据的对抗数据集，用于验证框架在实际场景中的适用性。该数据集包含96例去标识化的完整病历，涵盖病史、出院小结、出院诊断和出院处方四类核心临床信息。实验结果发现，真实临床数据中的非规范表述和噪声降低了系统的防御能力，在通用基准上表现尚可的模型在真实场景下面临更高的安全风险。人机回环验证机制显示，人工评估与自动化评估的Kappa一致性系数达到0.68-0.72，验证了框架评估结果的可靠性。

综上而言，本文构建了面向医疗多智能体系统的自动化安全评估框架MLSE，提出了级联错误传播的理论模型，并通过单体模型和多智能体系统两个层面的实验验证了医疗大模型面临的安全脆弱性。研究成果不仅揭示了医疗大模型在多智能体协作场景下面临的严峻安全挑战，也为医疗人工智能系统的安全部署提供了重要的理论依据和实践参考。

\section{主要创新点}
\label{sec:contributions}

本文的研究致力于解决当前医疗大语言模型安全评估领域中存在的问题：缺乏面向多智能体协作场景的安全评估框架、级联错误传播缺乏理论建模、以及评估结果缺乏真实临床数据验证。总体而言，本文工作的主要创新点如下：

（1）提出了面向医疗多智能体协作系统的自动化安全评估框架MLSE。区别于现有HarmBench、AgentHarm等框架主要关注单体模型或通用场景的安全评估，MLSE框架创新性地将医疗临床诊疗流程抽象为链式多智能体协作结构，实现了对真实临床场景的精准模拟。框架采用模块化设计，支持数据投毒、后门攻击、提示词注入和越狱攻击四类典型安全威胁的评估，并通过六类专业化智能体实现了评估流程的全自动化。实验表明，该框架能够有效识别医疗大模型在多智能体场景下的安全漏洞。

（2）首次提出了医疗协作系统中错误放大因子$\alpha$的数学定义及其理论模型。针对多智能体系统中级联错误传播的定量分析问题，定义了第$i$个智能体向第$i+1$个智能体传递错误时的放大因子$\alpha_i$。该因子衡量了在存在上游错误时下游智能体发生错误的条件概率与无上游错误时发生错误的条件概率之比，当$\alpha_i$大于1时表示错误在传递过程中被放大。16种配置组合的实验结果表明，在医疗多智能体协作系统中，错误放大因子显著大于1，验证了本文提出的放大假设。此外，本文还对传统评估指标进行了针对医疗场景的重新解释，将攻击成功率定义为安全下界代理指标，将困惑度解读为攻击隐蔽性指标，为医疗模型安全评估提供了新的理论视角。

（3）构建了首个基于真实临床数据的医疗多智能体对抗数据集。该数据集基于96例类风湿关节炎患者的去标识化病历构建，相比MedQuAD等通用医疗问答数据集，真实临床数据具有显著的噪声和非规范性，能够更准确地反映医疗大模型在实际应用场景中面临的安全挑战。案例研究揭示了诱导性用药推荐等长尾安全风险，为后续研究提供了重要的实证数据资源。

综上而言，本文在框架设计、理论建模和数据资源三个层面形成了较为完整的研究体系。MLSE框架为医疗大模型的安全性评估提供了一种通用、可扩展的解决方案；错误放大因子模型填补了多智能体系统安全评估的理论空白；真实临床对抗数据集则为该领域的后续研究奠定了数据基础。

\section{未来工作展望}
\label{sec:future}

本文分析了医疗多智能体协作系统安全评估领域中存在的问题，并针对问题提出相应解决方法，进行相关实验验证。随着医疗人工智能领域的不断发展，未来工作可以聚焦在以下几个方面：

（1）扩展数据规模与覆盖范围。本文构建的真实临床数据集包含96例类风湿关节炎样本，样本量和病种覆盖的有限性影响了研究结论的普适性。未来工作可与多家医疗机构开展合作，构建涵盖心血管疾病、肿瘤、神经系统疾病等多病种的大规模临床对抗数据集，提升研究结论的代表性。

（2）探索防御机制与安全增强技术。本文主要聚焦于安全漏洞的发现和评估，防御机制的研究相对薄弱。未来可研究智能体自我反思与自我修正技术、多智能体交叉验证机制、以及针对特定攻击类型的对抗性训练方法，形成攻防对抗的闭环。

（3）扩展到多模态医疗场景。实际临床诊疗涉及医学影像、检验报告波形数据等多模态信息。未来工作可研究医学影像注入攻击对多模态医疗大模型的影响，探索文本与图像之间的跨模态攻击传播机制，建立多模态综合安全评估框架。

（4）深化理论模型与推进工程应用。错误放大因子$\alpha$的数学模型仍相对简化，未考虑智能体之间的非线性相互作用和反馈机制。未来可引入随机过程、图论等数学工具建立更精确的理论框架。同时，推动MLSE框架向产业级安全评估平台演进，解决大规模并发测试、标准化报告等工程化问题。
